\scanpage{148}
Il est évident que la condition envisagée ne dépend pas du choix de l'immersion
$K\hookrightarrow\overline{\mathbb Q}$, \deleted{\unsure{donc}} i.e. du choix de $y$,
et peut donc s'écrire
\[
(*)\qquad(\tau\times1)\Sigma=\Sigma'\qquad\text{pour tte conjugaison }\tau
\]
\[\deleted{(\tau\times1)\Sigma\subset}\]
Je dis qu'elle équivaut aussi à
\[
(**)\qquad(1\times\tau)\Sigma\subset\Sigma'\qquad\text{pour tte conjugaison }\tau
\]
En effet, comme $(\tau\times\tau)\Sigma=\Sigma$, la relation $(*)$ équivaut à
$(\tau\times1)(\tau\times\tau)\Sigma=\Sigma'$ i.e. $(1\times\tau)\Sigma=\Sigma'$.

\uline{Donc la condition envisagée est symétrique en} $E,K$.

\par\noindent
\deleted{\mbox{\begin{minipage}[c]{\linewidth}
Considérons maintenant le sous-corps $K_0$ de $K$, corps de définition de
$\Sigma_y$: c'est \inserted{le corps des invariants du sous-groupe $H\subset G$,}
\unsure{par} des $g\in G$ tels que
\[
\deleted{g\Sigma_y\subset\Sigma_y}\qquad
(g\times1)\Sigma_y=\Sigma'_y,
\]
ce qui \deleted{implique} \inserted{équivaut}
\deleted{pour tte conjugaison} (comme $K_0$ donc $H$ doit être indépendant du choix de $y$) à
\[
\deleted{(g\times1)(\unsure{g^{-1}|g}\times g^{-1})\Sigma_y=\Sigma_y}\qquad
(g\times1)\Sigma\subset\unsure{\Sigma'|\Sigma}.
\]
\deleted{\unsure{par ci-dessus}\quad $\unsure{gK}$}
[vérification: $\Sigma_{y'}=(h\times h)\Sigma_y$ pour $h$ tel que $hy=y'$, donc
\[
(g\times1)\Sigma_{y'}=(g\times1)(h\times h)\Sigma_y=
\deleted{(g\times1)\Sigma_y}
(1\times h)(gh\times1)\Sigma_y=(1\times h)\Sigma'_y=\Sigma'_{y'},
\]
et comme $\Sigma=(g^{-1}\times g^{-1})\Sigma$, cela équivaut :
$(g\times1)(g^{-1}\times g^{-1})\Sigma=\Sigma'$,
i.e. à $(1\times g^{-1})\Sigma=\Sigma$, ou encore $(1\times g)\Sigma=\Sigma$]. Donc
\end{minipage}}}
\par

\uline{NB} Soit $\mathfrak U$ une structure CM sur le couple $E,K$,
et considérons $E\hookrightarrow E'$, $K\hookrightarrow K'$ des extensions de corps.
Alors l'idéal $\mathfrak U'=\mathfrak U(E'\otimes K')$ de $E'\otimes K'$ est une structure
\inserted{CM} de Sh. Tan. sur $E',K'$, et \uline{réciproquement}.
\deleted{La théorie de Galois}

\deleted{On dit que} Soit $\mathfrak U$
\deleted{structure CM sur $E,K$} \inserted{un idéal de $E\otimes K$}, envisagé
comme idéal de $E_K$; $\mathfrak U$ a un ``corps de définition'' $K_0\subset K$.
Je dis que celui-ci est invariant par extension du corps $E\longrightarrow E'$
(vérification triviale). De ceci
\margin{147}
